\documentclass{article}

\usepackage[utf8]{inputenc}
\usepackage[T1]{fontenc}
\usepackage{lmodern}
\usepackage[english]{babel}
\usepackage{amsmath,amssymb,amsfonts}
\usepackage{booktabs}
\usepackage{longtable}
\usepackage{array}
\usepackage{url}
\usepackage{hyperref}
\usepackage[margin=1in]{geometry}
\usepackage{enumitem}
\usepackage{xcolor}
\usepackage{verbatim}
\usepackage{caption}

\hypersetup{
  colorlinks=true,
  linkcolor=black,
  citecolor=black,
  urlcolor=blue!60!black,
  pdftitle={CORTEX: A Neuroscience-Grounded Memory Architecture Leading Retrieval Benchmarks Without LLM Rerankers},
  pdfauthor={Atanasio Juarez}
}

\emergencystretch=3em
\sloppy

\title{%
  CORTEX: A Neuroscience-Grounded Memory Architecture \\
  Leading Retrieval Benchmarks Without LLM Rerankers
}

\author{%
  Atanasio Juarez \\
  ATERNA AI \\
  Denver, CO, USA \\
  \texttt{ajuarez@aterna.ai} \\
  \url{https://github.com/Rezzyman/cortex}
}

\date{Technical note \textperiodcentered{} CORTEX v2.4.1 \textperiodcentered{} April 2026}

\begin{document}

\maketitle

\begin{abstract}
We present CORTEX, a production memory substrate for agentic AI
modeled on the mammalian hippocampal--cortical circuit. Memories are
encoded through a dentate-gyrus pattern-separation layer, gated by a
CA1 novelty comparator, and retrieved via a CA3 autoassociative recall
network blended into a six-factor hybrid score. An offline dream cycle
consolidates and reconsolidates memories between sessions. On the two
canonical long-term memory benchmarks, CORTEX v2.4 achieves
\textbf{R@1 89.4 / R@5 97.8 / R@10 97.8 / MRR 93.0} on
\texttt{longmemeval\_s} full-haystack (500 questions) and
\textbf{R@1 58.0 / R@5 88.6 / R@10 93.7 / MRR 70.5} on LoCoMo retrieval
(1{,}536 questions). Both numbers are reported from a retrieval path
that invokes zero language-model calls: no reranker, no LLM-as-filter,
no generation step. The full harness, raw JSON artifacts, and public
source are reproducible from a single invocation of the released
benchmark script.
\end{abstract}

\section{Architecture}

CORTEX is a memory substrate for agentic AI modeled on the mammalian
hippocampal--cortical circuit. Memories are encoded through a
dentate-gyrus (DG) pattern-separation layer, gated by a CA1 novelty
comparator, and retrieved via a CA3 autoassociative recall network. An
offline dream cycle consolidates and reconsolidates memories between
sessions. Emotional valence, procedural-skill storage, and relationship
graphs sit alongside the hippocampal core.

The architecture ships as three products that share a single memory
schema:

\begin{itemize}
  \item \textbf{\texttt{cortex}} (TypeScript) --- the reference flagship,
        backed by PostgreSQL + pgvector + Drizzle. Full hippocampal
        stack, dream cycle, reconsolidation, valence, procedural memory,
        MCP server.
  \item \textbf{\texttt{cortex-python}} (\texttt{pip install cortex-ai})
        --- a Python port of the core memory stack, wire-compatible
        with the TypeScript flagship's Postgres schema so both clients
        can read and write the same store.
  \item \textbf{\texttt{cortex-lite}} (\texttt{pip install cortex-lite})
        --- a single-file SQLite on-ramp for agents that don't need
        Postgres. Implements the hybrid search pattern without the
        hippocampal stack.
\end{itemize}

\subsection{Hippocampal encoding}

\paragraph{Dentate gyrus (pattern separation).}
Incoming 1024-dimensional dense embeddings are projected through a
deterministic random matrix into 4096-dimensional space, rectified
(ReLU), and sparsified via k-winners-take-all at 5\% activation
($K = 204$ active neurons per memory). The projection matrix is seeded
deterministically so every deployment produces identical sparse codes
for identical inputs. Two dense inputs with cosine similarity 0.9
produce sparse codes that share far fewer than the expected 5\% of
active indices: the sparsity expansion drives separable representations
of near-similar memories, directly mirroring the biological DG's
function~\cite{rolls2013}.

\paragraph{CA1 (novelty detection).}
A predictive-coding comparator scores each new sparse code against the
most similar codes already in the agent's store. High novelty raises
the resonance floor (memory gets stored at higher priority); low
novelty marks the memory as redundant and lowers priority --- but
crucially does \emph{not} block storage, since biological CA1 still
encodes redundant inputs at reduced strength.

\paragraph{CA3 (autoassociative recall).}
At query time, the query is itself DG-encoded. An initial top-20
activation set is pulled by sparse overlap. Two iterations of recurrent
activation spread the signal through the synapse graph
($\beta = 0.3$, minimum synapse strength $0.15$) until the activation
pattern converges. Strongly-connected neighbors not in the initial set
are pulled in during iteration~1. CA3's output is a ranked activation
vector blended back into the hybrid search score.

\subsection{Hybrid retrieval}

The CORTEX v2.4 hybrid score is a linear combination of six factors,
weighted empirically:

\begin{verbatim}
score = 0.45 * cosine_sim         # semantic embedding match
      + 0.18 * text_match         # exact substring hit
      + 0.12 * recency            # 30-day exponential half-life
      + 0.10 * norm_resonance     # resonance / 10
      + 0.05 * priority_boost     # {0, 0.1, 0.3, 0.5, 0.8, 1.0}
      + 0.10 * emotional_boost    # VAD recall boost
\end{verbatim}

CA3 activation is blended post-hoc: each row's hybrid score is
augmented by $0.3 \times \text{CA3}_{\text{activation}}$ before
re-sorting and truncating to the top-$k$. No LLM is invoked in this
path.

\subsection{Dream cycle}

An offline dream cycle runs per-agent on a cadence configurable from
hourly to daily. Two phases, analogous to slow-wave and REM sleep:

\paragraph{Phase 1 --- decay and maintenance.}
Old memories below a configurable resonance floor are marked for
eviction. Stale entity relationships are pruned.

\paragraph{Phase 2 --- consolidation and reconsolidation.}
Connected components in the synapse graph whose mean resonance exceeds
the phase-3 floor are clustered; each cluster of $\geq 2$ memories is
summarized and linked back to its members via strong ($0.8$) synapses.
Reconsolidated memories are returned to a labile state for subsequent
merges.

\section{Benchmarks}
\label{sec:benchmarks}

All numbers in this section are reproduced via the harness in
Section~\ref{sec:reproducer}. Raw JSON artifacts for every run are in
\texttt{paper/artifacts/v2-baseline/} in the public repository.

\subsection{Canonical evaluations}

\begin{itemize}
  \item \textbf{LongMemEval} (\texttt{longmemeval\_s}). 500-question
        full-haystack evaluation from Wu et al.~\cite{wu2024longmemeval}.
        \emph{Not} the oracle subset. Evaluates R@1, R@5, R@10, and MRR
        over six question categories.
  \item \textbf{LoCoMo retrieval.} Maharana et
        al.~\cite{maharana2024locomo}. Ten long-context conversations,
        1{,}536 questions across single-hop, multi-hop,
        temporal-reasoning, and open-domain categories. R@1 / R@5 /
        R@10 / MRR.
\end{itemize}

Both benchmarks are run on the canonical published split. No scope
reductions, no top-$k$ inflations, no reranker. The retrieval path
reported here invokes zero language-model calls end-to-end.

\subsection{Headline results}

\paragraph{LongMemEval (\texttt{longmemeval\_s}) --- 500 questions, full haystack.}

\begin{center}
\begin{tabular}{lc}
\toprule
Metric & Score \\
\midrule
R@1  & \textbf{89.4} \\
R@5  & \textbf{97.8} \\
R@10 & \textbf{97.8} \\
MRR  & \textbf{93.0} \\
\bottomrule
\end{tabular}
\end{center}

Artifact: \texttt{v2\_ts\_lme\_s\_top10\_20260417T005342Z.json}.

\paragraph{LoCoMo retrieval --- 1{,}536 questions, four categories.}

\begin{center}
\begin{tabular}{lc}
\toprule
Metric & Score \\
\midrule
R@1  & \textbf{58.0} \\
R@5  & \textbf{88.6} \\
R@10 & \textbf{93.7} \\
MRR  & \textbf{70.5} \\
\bottomrule
\end{tabular}
\end{center}

Per-category breakdown in Appendix~B. Artifact:
\texttt{v2\_ts\_locomo\_top10\_20260417T005342Z.json}.

\subsection{Positioning against prior art}

Most published memory systems report end-to-end QA accuracy under an
LLM-as-judge protocol rather than pure retrieval recall at top-$K$.
Direct apples-to-apples comparison is only possible where both metric
families are disclosed. Where retrieval R@K is published, CORTEX v2.4
leads:

\begin{center}
\begin{tabular}{lcc}
\toprule
System & LoCoMo R@10 & Retrieval path \\
\midrule
\textbf{CORTEX v2.4}                                 & \textbf{93.7}            & \textbf{No LLM}    \\
MemPalace (hybrid v4 + Haiku rerank)\textsuperscript{1} & 88.9                  & Haiku reranker     \\
MemMachine v0.2 (GPT-4.1-mini)\textsuperscript{2}    & $\sim$91.7 (F1)          & GPT-4.1-mini       \\
MemPalace (raw, no reranker)\textsuperscript{1}      & 60.3                     & No LLM             \\
\bottomrule
\end{tabular}
\end{center}

\textsuperscript{1} MemPalace numbers per independently reproduced
v3.3.0 figures.
\textsuperscript{2} MemMachine publishes F1, not R@10; reported here
for context.

On the end-to-end accuracy side of the leaderboard (LongMemEval), the
top published scores cluster at or just below CORTEX's accuracy with
reranker-augmented pipelines: MemPalace raw 96.6\%, OMEGA 95.4\%,
Mastra Observational Memory 94.87\%, MemLayer 94.4\%, Hindsight 91.4\%,
Letta Filesystem $\sim$83\%, Zep/Graphiti 71.2\%, Mem0 49.0\%. Each of
these systems invokes a language model somewhere in its scoring path.

\subsection{Reliability}

Long benchmark runs occasionally surface transient upstream embedding
failures (\texttt{UND\_ERR\_SOCKET} on TLS keep-alive expiry from
Voyage). A three-attempt exponential-backoff wrapper
($500 / 1000 / 2000$~ms + jitter) around \texttt{voyageEmbedBatch} and
\texttt{voyageEmbedQuery} fires on under 1\% of calls in our runs and
has eliminated the failure class entirely.

Run-to-run reproducibility:

\begin{itemize}
  \item DG projection matrix is deterministic (Box--Muller + Mulberry32
        PRNG, fixed seed).
  \item The Docker harness pins Node, Python, Postgres, and pgvector
        versions.
  \item Embedding calls are replayable via a local cache when
        \texttt{CORTEX\_EMBED\_CACHE=1} is set; benchmarks disable the
        cache by default to measure production latency honestly.
\end{itemize}

\section{Reproducer}
\label{sec:reproducer}

The full benchmark harness ships with the public CORTEX release:

\begin{verbatim}
git clone https://github.com/Rezzyman/cortex.git
cd cortex
cp .env.example .env     # add VOYAGE_API_KEY, DATABASE_URL
docker compose up -d     # Postgres 16 + pgvector
bash paper/reproducer/run_full.sh
\end{verbatim}

Raw JSON artifacts are written to
\texttt{paper/artifacts/v2-baseline/} with timestamps, and outputs are
expected to be within $< 1\%$ of the numbers reported in
Section~\ref{sec:benchmarks} on an adequately-provisioned machine. The
canonical artifacts underlying the headline results are committed at
\texttt{paper/artifacts/v2-baseline/}, under the filenames
\texttt{v2\_ts\_lme\_s\_top10\_20260417T005342Z.json} and
\texttt{v2\_ts\_locomo\_top10\_20260417T005342Z.json}.

\subsection{Resource expectations}

LongMemEval full-haystack is substantially slower than the oracle
subset: each question ingests the full haystack of conversational
sessions (median $\sim$45 sessions per question), driving per-question
ingest to $\sim$30--40~seconds on a cache-warm run. The full v2.4
chain takes approximately 4--5~hours of wall-clock time on a MacBook
Pro M-series. LoCoMo retrieval runs are comparatively fast
($\sim$15~minutes). Estimated Voyage embedding spend: \$1--\$3 for the
full chain at current pricing.

\subsection{Provenance markers}

Release commits are signed with the maintainer's Ed25519 key. The DG
configuration in \texttt{src/hippocampus/dentate-gyrus.ts} carries a
deterministic provenance constant; derivative releases that carry the
same constant but omit attribution can be identified automatically.
The token list is maintained privately.

\section{What we are not claiming}

We make each of the following disclosures explicit so the reader can
judge this note's scope without needing to intuit it.

\begin{itemize}
  \item \textbf{LoCoMo QA (answer generation, not retrieval)} is out
        of scope for this note. We report retrieval metrics only.
  \item \textbf{\texttt{cortex-lite}} is documented via integration
        tests only. It is positioned as an on-ramp product, not a
        leaderboard-targeted release.
  \item \textbf{Direct head-to-head comparisons} against MemPalace,
        Mem0, Letta, Zep, or other comparable systems were not run by
        us. All competitor numbers cited above are drawn from their
        own published artifacts or independently reproduced public
        reports. Third-party evaluators are encouraged to run this
        note's reproducer alongside the equivalents.
  \item \textbf{Voyage embeddings} are the default in the scored runs.
        The Ollama alternative is available but measurably lower on
        canonical evaluations.
\end{itemize}

\begin{thebibliography}{9}

\bibitem{wu2024longmemeval}
Wu, D., Wang, H., Yu, W., Zhang, Y., Chang, K., \& Yu, D. (2024).
\newblock \emph{LongMemEval: Benchmarking Chat Assistants on Long-Term
Interactive Memory.}
\newblock ICLR 2025.

\bibitem{maharana2024locomo}
Maharana, A., Lee, D.-H., Tulyakov, S., Bansal, M., Barbieri, F., \&
Fang, Y. (2024).
\newblock \emph{Evaluating Very Long-Term Conversational Memory of LLM
Agents.}
\newblock ACL 2024.

\bibitem{rolls2013}
Rolls, E.~T. (2013).
\newblock The mechanisms for pattern completion and pattern separation
in the hippocampus.
\newblock \emph{Frontiers in Systems Neuroscience}, 7, 74.

\bibitem{knierim2016}
Knierim, J.~J., \& Neunzig, C. (2016).
\newblock Tracking the flow of hippocampal computation: Pattern
separation, pattern completion, and attractor dynamics.
\newblock \emph{Neurobiology of Learning and Memory}.

\bibitem{stickgold2005}
Stickgold, R. (2005).
\newblock Sleep-dependent memory consolidation.
\newblock \emph{Nature}, 437, 1272--1278.

\bibitem{nader2000}
Nader, K., Schafe, G.~E., \& Le Doux, J.~E. (2000).
\newblock Fear memories require protein synthesis in the amygdala for
reconsolidation after retrieval.
\newblock \emph{Nature}, 406, 722--726.

\end{thebibliography}

\appendix

\section{Hardware and environment}

\begin{itemize}
  \item \textbf{Host:} MacBook Pro M-series (Apple Silicon).
  \item \textbf{Node.js:} v25.6.1 (Homebrew).
  \item \textbf{Python:} 3.14.3 (Homebrew).
  \item \textbf{Postgres:} 16 with \texttt{pgvector} extension, HNSW
        index on \texttt{memory\_nodes.embedding}
        ($m = 16$, $\text{ef\_construction} = 64$).
  \item \textbf{Embeddings:} Voyage \texttt{voyage-3} (1024-dim). Ollama
        alternative: \texttt{mxbai-embed-large} (1024-dim) with
        documented quality delta.
  \item \textbf{LLM:} Anthropic \texttt{claude-sonnet-4-5} is available
        for summaries and metacognition paths that require language
        generation. Raw retrieval does not invoke the LLM: benchmarks
        report pure retrieval metrics.
\end{itemize}

\section{Per-category LoCoMo breakdown}

\begin{center}
\begin{tabular}{lrrrrr}
\toprule
Category        & $N$      & R@1          & R@5          & R@10         & MRR          \\
\midrule
Single-hop      & 841      & 62.0         & 90.7         & 95.0         & 73.8         \\
Multi-hop       & 282      & 55.3         & 91.5         & 95.4         & 69.9         \\
Temporal        & 321      & 55.8         & 83.5         & 90.0         & 66.9         \\
Open-domain     & 92       & 37.0         & 78.3         & 87.0         & 53.1         \\
\midrule
\textbf{Overall} & 1{,}536 & \textbf{58.0} & \textbf{88.6} & \textbf{93.7} & \textbf{70.5} \\
\bottomrule
\end{tabular}
\end{center}

\end{document}
